\documentclass[11pt]{beamer}
\usetheme{Madrid}
\usepackage[utf8]{inputenc}
\usepackage{geometry}
\usepackage{hyperref}
\usepackage{amsmath, amssymb}
\usepackage{xcolor}

% Metadata
\title[Week 1]{Prompt Engineering: Week 1\\Introduction to Prompt Engineering}
\author{Mehmet Ali Akyol, PhD}
\institute{Ankara Medipol University}
\date{\\}

\begin{document}

% Title
\begin{frame}
  \titlepage
\end{frame}

% Overview
\begin{frame}
  \frametitle{Today: What We'll Explore}
  \begin{itemize}
    \item What is a "prompt"? A friendly, everyday definition
    \item Why prompts matter for everyone (not just programmers)
    \item Core terms: Prompt, LLM, Generative AI, Context
    \item Examples across fields: creative writing, law, marketing, engineering, education, medicine
    \item In-class activity: vague vs detailed prompts
    \item Common pitfalls and quick wins
    \item Assignment for next week
  \end{itemize}
\end{frame}

% What is a prompt?
\section{What is a Prompt?}
\begin{frame}
  \frametitle{What is a "Prompt"?}
  \begin{itemize}
    \item A \textbf{prompt} is how you talk to an AI system. It's the instruction, question, or example you give it.
    \item Everyday analogy: asking Google vs asking a friend. The way you ask changes the answer you get.
    \item Think of a prompt as a \textbf{brief} you give to a creative assistant.
    \item Good prompts provide: \textbf{purpose}, \textbf{context}, and \textbf{constraints}.
  \end{itemize}
\end{frame}

\begin{frame}
  \frametitle{Why Prompting Matters in 2025}
  \begin{itemize}
    \item AI is a general-purpose tool: writing, summarizing, analyzing, creating.
    \item Prompting is a new form of \textbf{digital literacy}.
    \item Every discipline benefits:
    \begin{itemize}
      \item Medicine: summarizing patient histories, drafting patient education.
      \item Law: drafting clauses, summarizing cases, spotting issues.
      \item Business: market research, pitch drafts, customer support.
      \item Engineering: code generation, documentation, debugging help.
      \item Design: brainstorming, mood boards, content ideas.
    \end{itemize}
    \item With better prompting, you get \textbf{clearer, safer, more useful} outputs.
  \end{itemize}
\end{frame}

% Core terms
\section{Core Terms}
\begin{frame}
  \frametitle{Key Terms (Plain English)}
  \begin{itemize}
    \item \textbf{Prompt}: The input message you give the AI.
    \item \textbf{LLM (Large Language Model)}: A system trained on lots of text to predict the next word.
    \item \textbf{Generative AI}: AI that can create new content (text, images, code) based on patterns it learned.
    \item \textbf{Context}: The information the AI considers for a response, including your conversation history.
  \end{itemize}
\end{frame}

% Examples
\section{Examples Across Fields}
\begin{frame}
  \frametitle{Example: Creative Writing}
  \begin{block}{Vague}
  Write a story about a journey.
  \end{block}
  \begin{block}{Better}
  You are a literary author. Write a 500-word short story in the style of magical realism about a paramedic in Istanbul who discovers time pauses whenever it rains. Use vivid sensory detail and a hopeful tone.
  \end{block}
  \begin{itemize}
    \item Why it's better: role, style, length, setting, tone, constraint.
  \end{itemize}
\end{frame}

\begin{frame}
  \frametitle{Example: Legal Drafting}
  \begin{block}{Vague}
  Draft a contract.
  \end{block}
  \begin{block}{Better}
  You are a contract lawyer. Draft a one-page independent contractor agreement for a freelance graphic designer in Turkey. Include scope of work, payment schedule (milestones), IP ownership (work-for-hire), confidentiality, and termination clauses. Use clear, plain language.
  \end{block}
  \begin{itemize}
    \item Why it's better: role, jurisdiction, length, required clauses, language.
  \end{itemize}
\end{frame}

\begin{frame}
  \frametitle{Example: Marketing}
  \begin{block}{Vague}
  Create a slogan for a coffee shop.
  \end{block}
  \begin{block}{Better}
  You are a brand strategist. Generate 10 short, punchy slogans for a student-friendly specialty coffee shop near a university in Ankara. Emphasize affordability, community, and late-night study vibes. Return as a numbered list.
  \end{block}
  \begin{itemize}
    \item Why it's better: role, audience, location, values, format.
  \end{itemize}
\end{frame}

\begin{frame}
  \frametitle{Example: Business}
  \begin{block}{Vague}
  Write a business plan.
  \end{block}
  \begin{block}{Better}
  You are a startup advisor. Draft a one-page business brief for a campus coffee cart at Ankara Medipol University. Include sections: problem, solution, target customers, value proposition, pricing, go-to-market (2 channels), basic cost estimate, and key risks with mitigations. Use bullet points.
  \end{block}
  \begin{itemize}
    \item Why it's better: scope and structure, audience, length, concrete sections, risk awareness.
  \end{itemize}
\end{frame}

\begin{frame}
  \frametitle{Example: Data Analysis}
  \begin{block}{Vague}
  Analyze this sales data.
  \end{block}
  \begin{block}{Better}
  You are a data analyst. Given a CSV with columns \texttt{date}, \texttt{product}, \texttt{region}, and \texttt{revenue}, compute monthly revenue by region for the last 12 months, identify any months more than 2 standard deviations from the 12-month regional mean, and return: (1) a 4-row table (one per region) with monthly totals, (2) a brief 80--120 word summary of anomalies, and (3) recommendations in 3 bullets.
  \end{block}
  \begin{itemize}
    \item Why it's better: clear inputs, time window, statistical rule, deliverables, and output format.
  \end{itemize}
\end{frame}

\begin{frame}
  \frametitle{Example: Engineering}
  \begin{block}{Vague}
  Write Python code to process data.
  \end{block}
  \begin{block}{Better}
  You are a software engineer. Write a Python function that reads a CSV file with columns \texttt{time} and \texttt{voltage}; compute a rolling average over a 5-second window; handle missing values by forward fill; include a docstring and type hints; return a pandas DataFrame; and show a short usage example.
  \end{block}
  \begin{itemize}
    \item Why it's better: role, file format, library, requirements, error handling, output spec, example.
  \end{itemize}
\end{frame}

\begin{frame}
  \frametitle{Example: Education}
  \begin{block}{Vague}
  Explain photosynthesis.
  \end{block}
  \begin{block}{Better}
  You are a high-school biology teacher. Explain photosynthesis for 9th graders in simple language using exactly 3 bullet points and one everyday analogy. Keep it under 120 words. End with two multiple-choice quiz questions (A--D) and the answer key.
  \end{block}
  \begin{itemize}
    \item Why it's better: role, audience level, structure, length, checks for understanding.
  \end{itemize}
\end{frame}

\begin{frame}
  \frametitle{Example: Medicine}
  \begin{block}{Vague}
  Summarize this patient case.
  \end{block}
  \begin{block}{Better}
  You are a clinical assistant. Given a patient note, produce a concise summary with sections: HPI, PMH, meds, allergies, vitals, key labs/imaging, assessment, and suggested next steps. Use neutral, non-diagnostic language; do not fabricate missing details; add a one-line disclaimer that this is not medical advice.
  \end{block}
  \begin{itemize}
    \item Why it's better: role, structured output, safety constraints, scope limits, disclaimer.
  \end{itemize}
\end{frame}

% Principles
\section{Principles of Good Prompts}
\begin{frame}
  \frametitle{Core Principles}
  \begin{itemize}
    \item \textbf{Clarity}: Make the goal explicit. What do you want?
    \item \textbf{Specificity}: Provide details: role, audience, length, style, format.
    \item \textbf{Context}: Give background, constraints, and examples.
    \item \textbf{Format}: Ask for structured output (bullets, table, JSON) when useful.
    \item \textbf{Tone}: Specify formal, friendly, persuasive, etc.
  \end{itemize}
\end{frame}

\begin{frame}
  \frametitle{Common Mistakes and Fixes}
  \begin{itemize}
    \item Too vague \textrightarrow{} add role, audience, and constraints
    \item Asking for everything at once \textrightarrow{} break into steps
    \item No examples \textrightarrow{} add 1--2 samples to guide style
    \item No format \textrightarrow{} request bullets or a template
    \item Ignoring limitations \textrightarrow{} ask for sources or confidence notes when needed
  \end{itemize}
\end{frame}

% Activity
\section{In-Class Activity}
\begin{frame}
  \frametitle{Activity: Vague vs Detailed Prompt}
  \begin{block}{How we'll run it}
    \begin{itemize}
      \item Pair up. Pick one task (e.g., summarize a news article, write an email, explain a concept).
      \item Write a \textbf{vague} prompt and run it.
      \item Improve it with role, audience, context, constraints, and format. Run again.
      \item Compare the outputs. What improved? What surprised you?
    \end{itemize}
  \end{block}
  \begin{block}{Debrief}
    Share one improvement that made the biggest difference.
  \end{block}
\end{frame}

% Discussion
\section{Discussion}
\begin{frame}
  \frametitle{Discussion Questions}
  \begin{itemize}
    \item Where in your field could prompting save time or improve quality?
    \item What risks or ethical concerns should we watch out for?
    \item How might you evaluate whether an AI output is "good enough" for your use?
  \end{itemize}
\end{frame}

% Key Takeaways
\section{Key Takeaways}
\begin{frame}
  \frametitle{Key Takeaways}
  \begin{itemize}
    \item Prompts are briefs: goal, context, constraints.
    \item Better prompts \textbf{consistently} lead to better outputs.
    \item Anyone can learn this skill with practice and iteration.
    \item Start simple, then refine with feedback from the output.
  \end{itemize}
\end{frame}

% Assignment
\section{Assignment}
\begin{frame}
  \frametitle{Assignment for Next Week}
  \begin{block}{Design 3 Prompts for the Same Task}
    \textbf{Task suggestion:} Summarize a news article or explain a concept from your field.
    \begin{itemize}
      \item Prompt 1: your first attempt (vague is fine)
      \item Prompt 2: refined with role, audience, constraints
      \item Prompt 3: add examples or a specific output format
    \end{itemize}
    Submit the 3 prompts and the 3 outputs as a single PDF with a short reflection (150--250 words) on what changed and why.
  \end{block}
\end{frame}

\begin{frame}
  \frametitle{Assignment Rubric (100\%)}
  \begin{itemize}
    \item \textbf{Clarity of task and prompts (25\%):} Goals are explicit; instructions are unambiguous.
    \item \textbf{Use of prompting levers (30\%):} Role, audience, context, constraints, and format are used effectively.
    \item \textbf{Iteration quality and reflection (25\%):} Clear progression from Prompt 1 to 3; reflection explains what improved and why.
    \item \textbf{Professionalism and formatting (10\%):} Clean structure; readable; outputs and prompts labeled.
    \item \textbf{Originality and appropriateness (10\%):} Creative yet suitable for the chosen task and audience.
  \end{itemize}
\end{frame}

% Thank You Slide
\begin{frame}{Questions and Discussion}
    \begin{center}
    {\Huge Thank You!}
    \vspace{1cm}

      \textbf{Dr.\\ Mehmet Ali Akyol}\\
      \texttt{mehmetali.akyol@gmail.com}
    \vspace{1cm}
    
    {\large Questions?}
    \end{center}
\end{frame}

\end{document}
